\section{The logical form of the intended theorem}
\label{sec:logic}

With the fixed predicates \eqref{eq:pos}--\eqref{eq:neg}, define
\begin{align*}
 \Rstr&:\ \forall n>5040\ \exists k\ \Pos(n,k),\\
 \Rdn&:\ \forall n>5040\ \neg\neg\exists k\ \Pos(n,k),\\
 \Rweak&:\ \forall n>5040\ \forall k\ \neg\Neg(n,k).
\end{align*}
The first is the intended positive strict Robin statement. It is
$\Pi^0_2$ in the language with fixed primitive-recursive function
symbols. The last is the absence of a certified strict violation,
a $\Pi^0_1$ formula after pairing its universal variables.

\begin{proposition}\label{prop:logic}
Intuitionistically, $\Rstr\to\Rdn\to\Rweak$.
The weak condition does not imply either of the other two in general.
\end{proposition}
\begin{proof}
The first implication introduces double negation. For the second,
suppose a negative certificate at $n$ is given. Incompatibility of
interval signs refutes every positive certificate at the same $n$,
contradicting $\neg\neg\exists k\,\Pos(n,k)$.

For the failure of a generic reverse rule, use the constant exact
interval $[0,0]$. It has no negative certificate and no positive
certificate, even under double negation. This example concerns a
logical rule; it makes no assertion that an actual Robin gap is zero.
\end{proof}

Replacing the strict inequality by $\Rweak$ thus changes the immediate
constructive content. Negative translation of a classical proof with
$\Rstr$ as conclusion ordinarily yields $\Rdn$, not $\Rstr$.
Removing the remaining double negation would be a Markov-type
witness-extraction step unless a bound or a separate theorem is supplied.
For example, encode a decidable predicate by the rapid rational name
which is zero until its first witness $j$ is found and then equals
$2^{-j}$. Positive apartness of its limit is precisely existence of a
witness. A generic double-negation removal rule would yield Markov's
principle for that predicate.
There is no need to speculate about the irrationality of $e^\gamma$,
and no such speculation would supply the missing proof.

\subsection{A concrete negative RH coding: the eta function}
\label{sec:etacert}

To avoid defining RH through Robin's criterion itself, use certificates
for zeros in rational rectangles strictly to the right of the critical
line. The following specifies a fixed primitive-recursive checker
design. Its analytic adequacy is a separate proof obligation, not a
completed result of the supplied Lean shell.

For $\Re s>0$ use the paired Dirichlet eta series
\[
 \eta(s)=\sum_{j=1}^{\infty}
       \bigl((2j-1)^{-s}-(2j)^{-s}\bigr).
\]
Let $R=[a,b]+i[c,d]$ be rational, with $1/2<a<b<1$, and let the
integer $M\ge1$ bound $|s|$ on $R$, for instance by the sum of the
largest absolute real and imaginary coordinates, rounded upward.
The identity
\[
 (2j-1)^{-s}-(2j)^{-s}
   =s\int_{2j-1}^{2j}t^{-s-1}\,dt
\]
yields the explicit tail
\begin{equation}\label{eq:eta-tail}
 \left|\eta(s)-\sum_{j=1}^{N}
       ((2j-1)^{-s}-(2j)^{-s})\right|
 \le\frac{M}{a}(2N+1)^{-a}
 \le\frac{2M}{\sqrt{2N}}.
\end{equation}
Indeed the omitted disjoint intervals are contained in
$[2N+1,\infty)$, and integration of $t^{-a-1}$ gives the first
inequality. Differentiating the displayed integrand gives
$(1-s\log t)t^{-s-1}$; the integrable majorant proves
\begin{equation}\label{eq:eta-derivative}
 |\eta'(s)|\le \frac1a+\frac{M}{a^2}\le 2+4M.
\end{equation}
These are explicit analytic estimates, with their finite-tail
counterparts available by the same inequalities.

For requested error $\varepsilon=2^{-k}$, choose
$N\ge8M^2\varepsilon^{-2}$. The tail is at most $\varepsilon/2$.
Evaluate each of the $2N$ finite terms with error at most
$\varepsilon/(4N)$. Each term is
$\exp(-s\log j)$, so the rational logarithm algorithm and a bounded
complex exponential Taylor sum suffice. For $|z|\le D$ with integer
$D\ge1$, a particularly crude adequate Taylor cutoff is
\[
 K=2\max(2D^2,p+4).
\]
The factorial lower bound from the final half of the factors and a
geometric ratio bound show that the tail after degree $K$ is at most
$2^{-p-3}$. Here $D$ must bound the approximate argument as well.
For a fully specified, deliberately generous schedule set
\[
 \begin{split}
 N&=8M^2\,2^{2k},\qquad p=k+\ell(N)+4,\\
 D&=M(2N+1)+1,\qquad
 h=p+\ell(M)+2(D+1)+4 .
 \end{split}
\]
For each $j\le2N$, choose the lower endpoint of
$\operatorname{LogI}(j,h)$, use it in $-s\log j$, and evaluate the
complex Taylor polynomial through degree $K$ as above. Since
$\log j\le j$ and the approximation error is at most one, $D$
bounds both arguments. The input perturbation is at most
$3^{D+1}M2^{-h}<2^{-p-4}$, and the Taylor tail is at most
$2^{-p-3}$. Each term therefore has error below
$2^{-p}<\varepsilon/(4N)$. These large but fixed
primitive-recursive bounds specify the rational eta observations
without invoking a universal program evaluator.

\Needspace{6\baselineskip}
The certificate code contains:
\begin{enumerate}[leftmargin=*]
\item The rational rectangle $R$ and an ordered rational subdivision
$z_0,\ldots,z_t=z_0$ of its positively oriented boundary.
\item A precision $k$ and the canonical rational eta approximations
$w_j$ at these vertices, each with coordinate error at most
$\varepsilon$. The checker recomputes them; it does not trust them.
\item A rational ray direction avoiding the polygon vertices, used to
compute an exact winding number.
\end{enumerate}
Put $L=2+4M$ and
\[
 \delta_j=|\Re(z_{j+1}-z_j)|+|\Im(z_{j+1}-z_j)|,\qquad
 E_j=2\varepsilon+L\delta_j.
\]
Require, for each $j$, that the coordinate square centered at $w_j$
of radius $E_j$ avoid zero. This is the rational disjunction
$|\Re w_j|>E_j$ or $|\Im w_j|>E_j$.
The eta image of the boundary segment and the straight polygon chord
from $w_j$ to $w_{j+1}$ both lie in that square. A straight homotopy
therefore avoids zero, and the endpoint homotopies agree between
adjacent segments.

Multiply all values by the complex conjugate of the chosen nonzero
Gaussian-rational direction. This is an orientation-preserving
rational similarity taking that direction to the positive real ray.
Require every transformed vertex to have nonzero imaginary coordinate.
Polygon winding is then computed by signed crossings.
For an edge with endpoint imaginary parts of opposite signs, its
real coordinate at height zero is
\[
 \frac{x_jy_{j+1}-x_{j+1}y_j}{y_{j+1}-y_j}.
\]
An upward crossing on the positive ray contributes $+1$; a downward
one contributes $-1$. Degenerate cases are rejected or avoided by
the supplied direction. Every decision is rational arithmetic.
The checker accepts when all the checks pass and this integer
winding number is positive. Malformed codes return false.
Call this fixed predicate $\Bad(c)$.

The analytic argument principle would establish soundness:
an accepted code encloses a zero of $\eta$. On $b<1$,
\[
 \eta(s)=(1-2^{1-s})\zeta(s),\qquad
 |2^{1-s}|\ge2^{1-b}>1,
\]
so the factor is separated from zero and the enclosed zero is a
zeta zero to the right of the critical line.

Define
\begin{equation}\label{eq:rhc}
 \RHc:\Longleftrightarrow \forall c\ \neg\Bad(c).
\end{equation}
This is a negative $\Pi^0_1$ formula and is intuitionistically stable:
$\neg\neg\RHc\to\RHc$. To see this, fix $c$ with $\Bad(c)$; it
would refute $\RHc$, contrary to the given double negation.

\paragraph{Adequacy that must still be proved.}
The finite evaluator, the homotopy argument, and the argument principle
need certified coded versions. Completeness requires surrounding every
off-line zero with a rational rectangle whose boundary has a positive
lower modulus, then finding an adequate finite subdivision. The functional
equation reflects left-half-strip zeros to the right. If the target RH
definition ranges beyond zeros in $0<\Re s<1$, the zero-free boundary
and exterior, and the trivial zeros, also need the usual interfaces.
None of these internal derivations is supplied by giving the checker
a name. No completed implementation of this checker is claimed here.

In particular the positive assertion that every zero-free compact
rectangle has a finite positive lower-bound certificate must not be
silently substituted for the negative assertion \eqref{eq:rhc}.
That conversion has its own proof requirements.

\subsection{Exactly sufficient certificate transformations}

Suppose four fixed primitive-recursive functions $u,p,v,q$ have been
constructed and $\ipra$ proves
\begin{align}
 n>5040&\ \to\
       \bigl(\Bad(u(n))\vee\Pos(n,p(n))\bigr),\label{eq:F}\\
 \Bad(c)&\ \to\
       \bigl(v(c)>5040\ \wedge\ \Neg(v(c),q(c))\bigr).\label{eq:R}
\end{align}
These statements must concern the actual specified checkers. They are
not axioms adopted in this report.

\begin{theorem}[Conditional assembly]\label{thm:assembly}
If the indicated object proofs of \eqref{eq:F}, \eqref{eq:R}, and
sign incompatibility are supplied, then $\ipra$ proves
$\RHc\leftrightarrow\Rstr$.
\end{theorem}
\begin{proof}
Assume $\RHc$ and fix $n>5040$. It rules out the first disjunct of
\eqref{eq:F}; $p(n)$ is the required positive witness.
Conversely assume $\Rstr$ and a particular $\Bad(c)$.
Equation \eqref{eq:R} gives a negative certificate at $v(c)>5040$,
while $\Rstr$ gives a positive certificate at that same integer.
Their incompatibility gives false.
\end{proof}

The theorem is compiled in the attached Lean file with both
transformations as explicit hypotheses. It is useful proof assembly,
not Robin's theorem with the analytic work concealed.

There is also a less witness-intensive legitimate converse. Because
$\RHc$ is stable, an intuitionistic proof of
$\neg\RHc\to\neg\Rstr$ suffices for $\Rstr\to\RHc$.
One must still prove that contrapositive in the correct theory;
it is not legitimate to turn $\neg\RHc$ into a positive
$\exists c\,\Bad(c)$ by an unstated Markov principle.

\subsection{A positive transfer metatheorem}

\begin{theorem}[Classical PRA is a sufficient intermediate target]
\label{thm:transfer}
Let $A\equiv\forall c\,a(c)$ and
$B\equiv\forall n\,\exists k\,b(n,k)$, where $a,b$ are
quantifier-free formulas over a fixed primitive-recursive language.
If classical first-order $\pra$ proves $A\leftrightarrow B$, then
$\ipra$ proves $A\leftrightarrow B$.
\end{theorem}
\begin{proof}
Classically, $A\to B$ implies the $\Pi^0_2$ sentence
\[
 \forall n\ \exists c,k\quad
       \bigl(\neg a(c)\vee b(n,k)\bigr).
\]
Herbrand witnessing for PRA supplies primitive-recursive terms
$c(n),k(n)$ and a quantifier-free PRA proof of
$\neg a(c(n))\vee b(n,k(n))$; a pair code handles the two witnesses.
The witnessing theorem is stated explicitly in
\cite[p.~13]{AvigadMining}. Quantifier-free arithmetic predicates
are decidable in $\ipra$, so this quantifier-free derivation can be
read intuitionistically. Under $A$ the first disjunct is impossible,
and $k(n)$ supplies the desired witness. This proves $A\to B$.

Apply the standard negative translation to the classical derivation
of $B\to A$. Decidability of quantifier-free formulas identifies
the translation of $A$ with $A$, and the translation of $B$ with
\[
 B^N\equiv\forall n\,\neg\neg\exists k\,b(n,k).
\]
Translated quantifier-free induction instances remain equivalent
to the permitted induction instances. Hence $\ipra$ proves
$B^N\to A$. Since $B\to B^N$ is intuitionistically valid, it also
proves $B\to A$.
\end{proof}

For Robin, fold the decidable condition $n>5040$ into $b(n,k)$ and
use a dummy witness below the cutoff. The theorem then applies to
$\RHc$ and $\Rstr$.

This metatheorem resolves one logical concern: strict witnesses are
not intrinsically incompatible with an initial classical approach.
Its premise is nevertheless a proof in \emph{classical PRA}, not
an arbitrary classical Lean proof. Establishing the necessary
arithmetic strength of the analytic argument remains the central
uncompleted task.

Likewise, conservation of a stronger theory for $\Pi^0_2$ sentences
can be applied to concrete witness statements such as
\eqref{eq:F}--\eqref{eq:R}. It cannot automatically be applied to
an implication with a $\Pi^0_2$ antecedent merely by calling the
whole implication arithmetical.
